\documentclass[12pt]{article}

\usepackage[spanish]{babel}
\usepackage[utf8]{inputenc}
\usepackage[T1]{fontenc}
\usepackage{amsmath, amssymb, amsthm}
\usepackage{geometry}
\usepackage{enumitem}
\usepackage{xcolor}
\usepackage{tikz}
\usepackage{hyperref}
\hypersetup{colorlinks=true, linkcolor=blue, urlcolor=blue}

\geometry{margin=2.5cm}

\newtheorem{teorema}{Teorema}
\newtheorem{definicion}{Definición}
\newtheorem*{ejemplo}{Ejemplo}
\newtheorem*{tarea}{Tarea}

\title{\bfseries Teorema del Residuo\\ \large y su aplicación a funciones zeta espectrales}
\author{Notas de clase}
\date{}

\begin{document}
\maketitle

\tableofcontents
\vspace{1cm}

\section{Motivación}

En análisis complejo, cuando una función $f(z)$ es holomorfa salvo en un número finito de puntos aislados (sus \textbf{polos}), es posible calcular integrales de contorno cerradas sin necesidad de parametrizar el camino de integración: basta con conocer el comportamiento local de $f$ alrededor de cada polo. Esta idea es la base del \textbf{teorema del residuo}, una de las herramientas más poderosas del análisis complejo.

Más adelante en esta clase usaremos esta misma maquinaria para estudiar \textbf{funciones zeta espectrales}, que generalizan la función zeta de Riemann a los autovalores de un operador diferencial (como el Laplaciano).

\section{El teorema del residuo}

\begin{definicion}[Residuo]
Si $f(z)$ tiene un polo aislado en $z_0$, el \textbf{residuo} de $f$ en $z_0$, denotado $\mathrm{Res}(f,z_0)$, es el coeficiente $a_{-1}$ de la serie de Laurent de $f$ alrededor de $z_0$:
\[
f(z) = \sum_{n=-\infty}^{\infty} a_n (z-z_0)^n, \qquad \mathrm{Res}(f,z_0) := a_{-1}.
\]
\end{definicion}

\begin{teorema}[Teorema del residuo]
Sea $f(z)$ una función meromorfa en un dominio que contiene un contorno cerrado simple $\omega$ (orientado positivamente), con polos $z_0, z_1, \dots, z_n$ \textbf{en el interior} de $\omega$ y sin singularidades sobre $\omega$. Entonces:
\[
\oint_{\omega} f(z)\,dz \;=\; 2\pi i \sum_{k} \mathrm{Res}(f, z_k).
\]
\end{teorema}

\noindent \textbf{Idea:} en vez de integrar directamente, "sumamos'' la contribución de cada singularidad encerrada, multiplicada por $2\pi i$.

\subsection{Cómo calcular un residuo}

En la práctica no siempre queremos escribir la serie de Laurent completa. Hay fórmulas directas según el \textbf{orden del polo}.

\begin{itemize}[leftmargin=1.2cm]
    \item \textbf{Polo simple (orden 1) en $z_0$:}
    \[
    \mathrm{Res}(f, z_0) = \lim_{z \to z_0} (z - z_0)\, f(z).
    \]

    \item \textbf{Polo de orden $k$ en $z_0$:}
    \[
    \mathrm{Res}(f, z_0) = \frac{1}{(k-1)!} \lim_{z \to z_0} \frac{d^{\,k-1}}{dz^{\,k-1}} \Big[ (z - z_0)^k f(z) \Big].
    \]
\end{itemize}

La fórmula de polo simple es, de hecho, el caso $k=1$ de la fórmula general.

\section{Ejemplo resuelto}

\begin{ejemplo}
Calcular
\[
\oint_{C} \frac{e^z}{z(z-1)^2}\,dz,
\]
donde $C$ es un contorno cerrado que encierra a los dos polos de la función: $z=0$ y $z=1$.
\end{ejemplo}

\noindent \textbf{Paso 1. Identificar los polos.}
La función
\[
f(z) = \frac{e^z}{z(z-1)^2}
\]
tiene:
\begin{itemize}
    \item un polo \textbf{simple} en $z_0 = 0$ (el factor $z$ aparece una sola vez), y
    \item un polo de \textbf{orden 2} en $z_1 = 1$ (el factor $(z-1)^2$).
\end{itemize}

\noindent \textbf{Paso 2. Residuo en $z=0$ (polo simple).}
\[
\mathrm{Res}(f,0) = \lim_{z \to 0} z \cdot \frac{e^z}{z(z-1)^2} = \left.\frac{e^z}{(z-1)^2}\right|_{z=0} = \frac{e^0}{(0-1)^2} = \frac{1}{1} = 1.
\]

\noindent \textbf{Paso 3. Residuo en $z=1$ (polo de orden 2).}
Usamos la fórmula con $k=2$:
\[
\mathrm{Res}(f,1) = \frac{1}{(2-1)!}\lim_{z \to 1} \frac{d}{dz}\left[ (z-1)^2 \cdot \frac{e^z}{z(z-1)^2} \right] = \lim_{z\to 1} \frac{d}{dz}\left[\frac{e^z}{z}\right].
\]
Derivamos con la regla del cociente:
\[
\frac{d}{dz}\left[\frac{e^z}{z}\right] = \frac{e^z \cdot z - e^z \cdot 1}{z^2} = \frac{e^z (z-1)}{z^2}.
\]
Evaluamos el límite en $z=1$:
\[
\lim_{z\to 1} \frac{e^z(z-1)}{z^2} = \frac{e^1 \cdot (1-1)}{1^2} = \frac{e \cdot 0}{1} = 0.
\]
Por lo tanto $\mathrm{Res}(f,1) = 0$.

\noindent \textbf{Paso 4. Aplicar el teorema del residuo.}
\[
\oint_C \frac{e^z}{z(z-1)^2}\,dz = 2\pi i \big[\mathrm{Res}(f,0) + \mathrm{Res}(f,1)\big] = 2\pi i (1 + 0) = \boxed{2\pi i}.
\]

\section{La función zeta de Riemann: ecuación funcional}

Antes de definir la función zeta espectral, recordamos la función zeta de Riemann $\zeta(s)$, definida originalmente para $\mathrm{Re}(s) > 1$ por
\[
\zeta(s) = \sum_{n=1}^{\infty} \frac{1}{n^s},
\]
y luego extendida (por continuación analítica) a todo el plano complejo salvo un polo simple en $s=1$. Esta extensión satisface la célebre \textbf{ecuación funcional}:
\[
\zeta(s) = 2^{s}\, \pi^{s-1} \, \sin\!\left(\frac{\pi s}{2}\right) \, \Gamma(1-s)\, \zeta(1-s).
\]

Esta identidad relaciona el valor de $\zeta$ en $s$ con su valor en $1-s$, y es la que permite estudiar $\zeta(s)$ fuera de la región original de convergencia $\mathrm{Re}(s)>1$.

\section{Funciones zeta espectrales}

La misma idea de "sumar sobre un espectro'' se puede aplicar a operadores diferenciales, no solo a los números naturales.

\begin{definicion}[Espectro discreto]
Sea $\Delta$ un operador (por ejemplo, el Laplaciano) cuyo espectro es \textbf{discreto}, es decir, sus autovalores forman un conjunto numerable
\[
\Lambda = \{\lambda_n\}_{n=1}^{\infty}, \qquad 0 < |\lambda_1| \leq |\lambda_2| \leq |\lambda_3| \leq \cdots
\]
\end{definicion}

\begin{definicion}[Función zeta espectral]
Se define la función zeta asociada al espectro $\Lambda$ como
\[
\zeta_{\Delta}(s) = \sum_{n=1}^{\infty} \frac{1}{\lambda_n^{s}} = \sum_{n=1}^{\infty} \lambda_n^{-s}.
\]
\end{definicion}

Esta serie converge para $\mathrm{Re}(s)$ suficientemente grande (de forma análoga a $\zeta(s) = \sum n^{-s}$), y define una función holomorfa en esa región. La pregunta natural —planteada en clase— es si $\zeta_\Delta(s)$ admite también una extensión (continuación analítica) más allá de su región de convergencia, y si esa extensión debe interpretarse como un \textbf{valor principal}.

Se introdujo además la notación
\[
\xi = \sigma + i t
\]
para la variable compleja, y se señaló una \textbf{abscisa de convergencia} $a_c = 1$: el valor de $\mathrm{Re}(\xi)$ a partir del cual la serie que define a $\zeta_\Delta$ converge.

\section{Tarea}

\begin{tarea}
¿Cuál es el residuo (o el polo) de $\zeta_{\Lambda}(s)$?
\end{tarea}

\noindent \textit{Sugerencia:} recordar que $\zeta(s)$ de Riemann tiene un único polo, simple, en $s=1$, con residuo
\[
\mathrm{Res}(\zeta, 1) = 1.
\]
La tarea consiste en investigar si $\zeta_\Lambda(s)$ tiene un comportamiento análogo: ¿tiene un polo simple?, ¿en qué valor de $s$?, ¿cuál es su residuo, y qué información geométrica o espectral codifica ese residuo?

\end{document}